\uuid{OI45}
\exo7id{422}
\titre{exo7 422}
\auteur{ridde}
\organisation{exo7}
\datecreate{1999-11-01}
\isIndication{false}
\isCorrection{false}
\chapitre{Polynôme, fraction rationnelle}
\sousChapitre{Racine, décomposition en facteurs irréductibles}
\cpge{mpsi-mp2i.nombres-complexes.f, pcsi.nombres-complexes.f, ptsi.nombres-complexes.f, mpsi-mp2i.polynomes-fractions-rationnelles.c, pcsi.polynomes.c, ptsi.polynomes.c, mpsi-mp2i.polynomes-fractions-rationnelles.f, pcsi.polynomes.e, ptsi.polynomes.e}
\module{Algèbre}
\niveau{L1}
\difficulte{}

\contenu{
\texte{
Soit $n \in \Nn^*$ et $P (X) = (X + 1)^n- (X-1)^n$.
}
\begin{enumerate}
    \item \question{Quel est le degr\'e de $P$ ?}
    \item \question{Factoriser $P$ dans $\Cc[X]$.}
    \item \question{Montrer que $\forall p \in \Nn^* \, \, \prod\limits_{k = 1}^p
\text{cotan} (\dfrac{k \pi}{2p + 1}) = \dfrac 1{\sqrt{2p + 1}}$.}
\end{enumerate}
}
